\section{A finite perturbation method for exact minimal complexes}
\label{sec:perturbation-theory}

The window scanner reduces the portion of a complex that must be
constructed.  A complementary question is how to reduce a complex once it
exists.  This section separates scalar linear algebra from the remaining
cobordism calculation and derives a finite homological perturbation
algorithm.  The construction uses the established perturbation lemma
\cite{crainic}; the point here is an explicit nilpotence bound for the
implemented algebra, an exact multiplicity profile, and a checked
implementation of the resulting maps.

All objects below have the same current frontier with $w=2k$ endpoints.
They are loopless matchings, with raw homological placements retained and
the quantum grading forgotten.  The grading introduced next is an
intrinsic grading of the morphism algebra.  It is not the raw homological
degree and does not require the program to reconstruct the forgotten
quantum shifts.

\subsection{Positive grading without a Catalan frontier assumption}

For two matchings $a,b$, let $c(a,b)$ be the number of circle components in
the union of $a$ and a reversed copy of $b$.  The implemented morphism
basis consists of one canonical disk on each such circle, decorated with
zero or one dot.  This is the familiar dotted-cobordism and arc-algebra
description\cite{barnatan,khovanov_arc}.  A basis element carrying $d$
dots receives intrinsic degree
\begin{equation}
 \ell(a,b;d)=k-c(a,b)+2d.
 \label{eq:intrinsic-degree}
\end{equation}

The relevant matching family is the one occurring in the executed scan.
Its endpoints need not lie on one noose, and its matchings need not be
noncrossing with respect to their numerical labels.  Nevertheless the
spherical projection has a checkerboard face coloring.  Orient each edge
so that the black face is on its left.  At each crossing the four edge
directions alternate between incoming and outgoing, and either smoothing
joins incoming to outgoing.  At a fixed cut this gives a fixed division of
the frontier into $k$ positive and $k$ negative endpoints; every smoothing
matching joins opposite classes.

Thus the abstract disk gluings used below are oriented.  This is the
condition needed for the Euler-characteristic calculation, and it does
not require a Catalan count.  One should not instead declare every
unconstrained perfect matching on a labeled set admissible: for four
endpoints the three distinct pairings together need not admit the required
common endpoint orientations.  The present argument uses matchings
actually furnished by a spherical diagram.

\begin{lemma}[Homogeneous composition and degree range]
\label{lem:intrinsic-grading}
Every nonzero basis morphism has intrinsic degree in $[0,w]$.
Composition adds intrinsic degrees.  Degree zero consists precisely of
scalar identities between equal matchings.
\end{lemma}

\begin{proof}
Since $1\leq c(a,b)\leq k$ for a nonempty frontier and
$0\leq d\leq c(a,b)$,
\[
 0\leq k-c(a,b)+2d\leq k+c(a,b)\leq2k=w.
\]
At the empty frontier only the empty object and its scalar endomorphisms
remain, so the statement also holds for $w=0$.  Degree zero requires
$c(a,b)=k$ and $d=0$.  The union has $k$ two-edge circles exactly when
$a=b$, and the undotted canonical cobordism is then the identity.

For composition from $a$ through $b$ to $c$, abbreviate
$c_{ab}=c(a,b)$, $c_{bc}=c(b,c)$ and $c_{ac}=c(a,c)$.  Before reduction,
the composite surface has Euler characteristic
\begin{equation}
 \chi=c_{ab}+c_{bc}-k,
 \label{eq:composite-euler}
\end{equation}
because $k$ intervals of the middle matching are glued.  Let $d_{\rm in}$
be the total input dot count.  Replacing this surface by its canonical
output disks through neck cutting raises Euler characteristic by two for
each additional dot.  Accordingly every nonzero output monomial satisfies
\begin{equation}
 2d_{\rm out}=2d_{\rm in}+c_{ac}-\chi.
 \label{eq:dot-euler-balance}
\end{equation}
The factor two in this identity is essential.  A positive-genus component
contributes zero over $\mathbb F_2$ after handle reduction; double dots and
undotted spheres also vanish, while a dotted sphere evaluates to one.
These relations preserve the same degree balance whenever a term remains.
Therefore
\begin{align*}
 \ell_{\rm out}
 &=k-c_{ac}+2d_{\rm out}\\
 &=k-\chi+2d_{\rm in}\\
 &=2k-c_{ab}-c_{bc}+2d_{\rm in}
   =\ell_{\rm first}+\ell_{\rm second}.
\end{align*}
The calculation concerns abstract oriented surfaces, so numerical endpoint
order and the number of boundary components of the scanned region do not
enter it.
\end{proof}

Let $J$ be the ideal of positive intrinsic-degree morphisms.  For a finite
set $\mathcal M$ of distinct executed matchings, the basic algebra is
$A=\operatorname{End}(\bigoplus_{a\in\mathcal M}a)$.

\begin{corollary}[Uniform nilpotence]
\label{cor:radical-nilpotence}
One has
\[
 J^{w+1}=0,\qquad A/J\cong\prod_{a\in\mathcal M}\mathbb F_2.
\]
The same nilpotence exponent applies to matrices of arbitrary finite size
over this matching category.  In the additive closure, an object with
$m_a$ copies of each matching has endomorphism quotient
$\prod_a\operatorname{Mat}_{m_a}(\mathbb F_2)$, rather than a product of
one-dimensional fields.  The ideal $J$ is the radical of these
finite-dimensional algebras.
\end{corollary}

\begin{proof}
A product of $w+1$ positive-degree basis terms would have degree at least
$w+1$, where all morphism spaces vanish by
Lemma~\ref{lem:intrinsic-grading}.  Expand sums and matrix products into
composable paths to obtain the same conclusion for arbitrary matrix
sizes.  The degree-zero description gives the displayed quotients.  A
nilpotent ideal lies in the radical, while a semisimple quotient forces
the reverse containment here.
\end{proof}

\subsection{The scalar differential gives the exact remaining dimensions}

Let $(C,d)$ be a finite complex over this category.  Extract $d_0$ by
keeping only the undotted constant coefficient between equal matchings;
put $\delta=d+d_0$.  Every entry of $\delta$ lies in $J$.  Projection to
degree zero is compatible with composition, so $d^2=0$ implies
$d_0^2=0$.

For each matching $a$ and raw degree $h$, let $V_{a,h}$ be the vector space
of its copies in $C^h$.  The scalar differential forms independent binary
complexes
\[
 \cdots\longrightarrow V_{a,h-1}
 \xrightarrow{d_{0,a,h-1}}V_{a,h}
 \xrightarrow{d_{0,a,h}}V_{a,h+1}\longrightarrow\cdots.
\]
Define
\begin{equation}
 m_{a,h}=\dim V_{a,h}
          -\operatorname{rank}(d_{0,a,h-1})
          -\operatorname{rank}(d_{0,a,h}).
 \label{eq:scalar-minimal-profile}
\end{equation}

\begin{theorem}[Exact minimal multiplicity profile]
\label{thm:scalar-minimal-profile}
There exists a complex homotopy equivalent to $(C,d)$ whose differential
lies entirely in $J$ and which has exactly $m_{a,h}$ copies of matching
$a$ in degree $h$.  Every such minimal model has the same multiplicities.
Thus the number of remaining objects can be computed by binary linear
algebra before carrying out the nonlinear cobordism correction.
\end{theorem}

\begin{proof}[Construction of the scalar contraction]
In every binary complex choose decompositions
\[
 V_{a,h}=B_{a,h}\oplus H_{a,h}\oplus L_{a,h},
 \qquad
 B_{a,h}=\operatorname{im}d_{0,a,h-1},
\]
where $B_{a,h}\oplus H_{a,h}=\ker d_{0,a,h}$ and
$d_0:L_{a,h}\to B_{a,h+1}$ is an isomorphism.  Let
$H=\bigoplus_{a,h}a^{m_{a,h}}[h]$, using $[h]$ here only to record a raw
placement.  Let $i:H\to C$ include the chosen homology representatives,
$p:C\to H$ project along $B\oplus L$, and $h:C\to C$ take each boundary
back to its chosen preimage in $L$, vanishing on $H\oplus L$.

All these maps have scalar identity entries between equal matchings.  They
satisfy
\begin{equation}
 pi=1,\quad d_0i=0,\quad pd_0=0,\quad
 d_0h+hd_0=1+ip,\quad hi=0,\quad ph=0,\quad h^2=0.
 \label{eq:scalar-contraction-identities}
\end{equation}
The next theorem lifts this contraction to the full $d$ without changing
$H$ as an object.  Its resulting differential lies in $J$.  Uniqueness is
proved below.
\end{proof}

The program obtains $B$ from independent original image columns and $L$
from their corresponding source vectors.  It extends $B$ to a basis of
the kernel to choose $H$, and solves for coordinates in $B\oplus H\oplus L$.
This provides the actual maps in~\eqref{eq:scalar-contraction-identities},
not only the ranks in~\eqref{eq:scalar-minimal-profile}.

\subsection{A finite homological perturbation formula}

Composition is written in the conventional order: $fg$ means first $g$,
then $f$.  The code's composition helper takes its arguments in execution
order: the second map acts after the first.

\begin{theorem}[Finite lifted contraction]
\label{thm:finite-hpl}
Assume~\eqref{eq:scalar-contraction-identities}, let $d=d_0+\delta$ with
$d^2=0$, and suppose the entries of $\delta$ lie in $J$.  Define
\begin{align}
 F&=\sum_{j=0}^{w}(h\delta)^j=(1+h\delta)^{-1},&
 G&=\sum_{j=0}^{w}(\delta h)^j=(1+\delta h)^{-1},
 \label{eq:finite-series}\\
 i'&=Fi,& p'&=pG,& h'&=Fh=hG,\nonumber\\
 D&=p\delta Fi=pG\delta i.&&
 \label{eq:finite-hpl-maps}
\end{align}
These are finite sums, and $(H,D)$ is a minimal complex with a strong
deformation retraction from $(C,d)$.  Explicitly,
\begin{gather}
 D^2=0,\quad di'=i'D,\quad p'd=Dp',\quad p'i'=1,
 \label{eq:lifted-chain-identities}\\
 dh'+h'd=1+i'p',\qquad h'i'=0,\quad p'h'=0,\quad (h')^2=0.
 \label{eq:lifted-retraction-identities}
\end{gather}
The number of perturbation factors needed is at most $w$, independent of
the number of objects or the size of their scalar blocks.
\end{theorem}

\begin{proof}
The maps $i,p,h$ have intrinsic degree zero.  Therefore $h\delta$ and
$\delta h$ have entries in $J$, and
Corollary~\ref{cor:radical-nilpotence} proves the inverse identities in
\eqref{eq:finite-series}.  These are the standard perturbation formulas
\cite[Sections~2--3]{crainic}; a direct verification in characteristic two
is short enough to record.

Put $e=ip$ and
\[
 A=\delta F=G\delta.
\]
The finite series imply
\begin{equation}
 A=\delta+\delta hA=\delta+Ah\delta,
 \quad F=1+hA,\quad G=1+Ah.
 \label{eq:auxiliary-hpl-identities}
\end{equation}
There is also the identity
\begin{equation}
 d_0A+Ad_0=AeA.
 \label{eq:auxiliary-hpl-maurer}
\end{equation}
To verify it, multiply its left side by $1+\delta h$ on the left and
$1+h\delta$ on the right.  Using $AF^{-1}=\delta$ and
$G^{-1}A=\delta$ gives
\begin{align*}
 &(1+\delta h)(d_0A+Ad_0)(1+h\delta)\\
 &\quad=d_0\delta+\delta d_0
       +\delta(hd_0+d_0h)\delta\\
 &\quad=\delta^2+\delta(1+e)\delta=\delta e\delta.
\end{align*}
Here $d^2=0$ gives $d_0\delta+\delta d_0=\delta^2$ over
$\mathbb F_2$.  Multiplying by $G$ and $F$ proves
\eqref{eq:auxiliary-hpl-maurer}.

Now $D=pAi$, so
\[
 D^2=pAeAi=p(d_0A+Ad_0)i=0.
\]
Since $i'=i+hAi$, equations
\eqref{eq:scalar-contraction-identities} and
\eqref{eq:auxiliary-hpl-identities} give
\[
 di'=eAi+hd_0Ai=eAi+hAeAi=i'D.
\]
The same calculation on the other side gives $p'd=Dp'$.  Expanding
$p'i'=(p+pAh)(i+hAi)$ leaves only $pi=1$, by the three side conditions
$hi=ph=h^2=0$.

Finally $h'=h+hAh$.  Expand $dh'+h'd$ and use
$A=\delta+\delta hA=\delta+Ah\delta$ to cancel the terms containing a
single $\delta$.  What remains is
\begin{align*}
 dh'+h'd
 &=1+e+eAh+hAe+h(d_0A+Ad_0)h\\
 &=1+e+eAh+hAe+hAeAh=1+i'p'.
\end{align*}
Expanding $h'i'$, $p'h'$ and $(h')^2$ similarly makes every term contain
$hi$, $ph$ or $h^2$, so each vanishes.  Each term of $D=p\delta Fi$
contains a positive-degree factor $\delta$, hence $D$ lies in $J$ and is
minimal.
\end{proof}

\begin{lemma}[Uniqueness of minimal models]
\label{lem:minimal-uniqueness}
Two finite complexes with differentials in $J$ that are homotopy equivalent
are chain isomorphic.  Their multiplicities agree matching by matching and
raw degree by raw degree.  In particular, the multiplicities in
Theorem~\ref{thm:scalar-minimal-profile} agree with any exhaustive unit
cancellation result.
\end{lemma}

\begin{proof}
Reduce a homotopy equivalence and its inverse modulo $J$.  Both reduced
differentials are zero, so the homotopy equations say that the reduced
term maps are inverse in each degree.  Thus the binary multiplicity spaces
agree.  Each lifted composite differs from the identity by a matrix in
$J$.  Such a matrix is inverted by a finite geometric series of length
at most $w+1$.  Consequently each original term map is invertible; its
inverse commutes with the differential because the term map does.  This
gives a chain isomorphism.  The assertion concerns multiplicities and
isomorphism type, not a canonical displayed basis or identical matrices.
\end{proof}

\subsection{What the implementation checks}

The research module \texttt{perturbation.py} first constructs the scalar
contraction, then computes $i'$ by repeatedly applying $h\delta$ to its
current term.  It obtains $D$ from $p\delta i'$.  With certification enabled
it also constructs $p',h'$ and verifies all identities in
\eqref{eq:lifted-chain-identities}--\eqref{eq:lifted-retraction-identities},
together with the original $d^2=0$, over the same exact cobordism algebra
used by the scanner.  It checks the finite-series termination and the
absence of scalar units in the output differential.

The corrected maps are stronger evidence than an agreement of final ranks:
they give the explicit homotopy equivalence needed for continuing the scan.
Because planar extension preserves homotopies, an installed model can be
extended through subsequent crossings.  The test suite compares its
matching-and-degree multiplicities with exhaustive scalar cancellation at
intermediate prefixes, checks final homology independently, and confirms
that a corrupted contraction map is rejected.  An additional independent
grading audit examined all bipartite matching monomial products through
six endpoints, including the non-Catalan six-endpoint matching: all 3533
products and 1576 nonzero output monomials obeyed
Lemma~\ref{lem:intrinsic-grading}.

\subsection{Cost, useful regimes, and remaining obstruction}

Write $N$ for the input object count at one stage and $H=\sum_{a,h}m_{a,h}$
for the output count.  There are at most $2^{w/2}$ dot-mask coefficients in
one morphism.  Scalar contraction uses binary matrices with at most $N$
rows and columns and therefore has polynomial bit cost, with a standard
cubic bound sufficient here.  The inclusion series performs at most $w$
iterations on matrices of size at most $N$ by $H$.  A conservative bound
for the reduction without optional certificates is
\begin{equation}
 O\bigl(N^3+(w+1)N^2(H+1)\,2^{O(w)}\bigr).
 \label{eq:hpl-cost}
\end{equation}
Computing and verifying all contraction maps has the coarser bound
\[
 O\bigl((w+1)(N+1)^3\,2^{O(w)}\bigr).
\]
These estimates allow dense maps; sparsity can help but is not assumed.
They follow by counting ordinary matrix-composition paths and charging
$2^{O(w)}$ for a cobordism product and its bit operations.  Cache storage
and dictionary bookkeeping remain polynomial in $N$ times $2^{O(w)}$.

The attraction is most pronounced when the scalar differential has large
rank, so $H\ll N$.  In the extreme $H=0$, scalar acyclicity already implies
full contractibility: the corrected inclusion and reduced differential
have empty source.  If an explicit contracting homotopy is requested, its
construction still has a cost.  When scalar cancellation is sparse and
cheap, assembling a global binary decomposition and dense corrected maps
can instead add overhead.  Thus no pointwise or asymptotic speed advantage
over the existing scalar scanner is claimed.  The experiment section
reports the measured regimes, including the distinction between natural
knot complexes and synthetic algebraic blocks.

This method does not control $N$ or $H$ by frontier width alone.  It
computes the minimal multiplicity profile of an already represented
complex, and an explicit output must still list its $H$ objects.  Known
fixed-braid-width exponential examples therefore remain an obstruction to
a universal explicit-basis complexity theorem.  The finite perturbation
bound supplies a precise local reduction cost and a certificate; combining
it with window bounds or a separate structural compression theorem is the
remaining route to stronger global guarantees.
